% Fichier engendré par tools/generate-tex.py à partir de 04-grandeurs-et-mesures.
% Ne pas éditer à la main : tout le texte provient des commentaires du fichier Lean.
\documentclass[11pt,a4paper]{article}

% Compile aussi bien avec pdflatex qu'avec xelatex ou tectonic.
\usepackage{iftex}
\ifPDFTeX
  \usepackage[T1]{fontenc}
  \usepackage[utf8]{inputenc}
\else
  \usepackage{fontspec}
\fi
\usepackage[french]{babel}
\usepackage{amsmath,amssymb,amsthm}
\usepackage[margin=2.5cm]{geometry}
\usepackage[hidelinks]{hyperref}

\theoremstyle{definition}
\newtheorem{definition}{Définition}
\theoremstyle{plain}
\newtheorem{theoreme}{Théorème}
\newtheorem{lemme}[theoreme]{Lemme}

% Renvoi à la déclaration Lean, une ligne après la démonstration, aligné à droite.
\newcommand{\source}[2]{\par\smallskip\noindent\hfill{\footnotesize\href{#1}{\texttt{#2}}}\par\smallskip}

\title{Grandeurs et mesures}
\date{}

\begin{document}
\maketitle
% Les identifiants Lean en machine à écrire ne se coupent pas : on tolère des
% espacements plus lâches plutôt que des lignes qui débordent.
\sloppy

\section{GrandeursEtMesures}

\noindent
Collège \textemdash{} section \og{} Grandeurs et mesures \fg{}.
Les formules d'aire et de volume sont ici des définitions, comme au collège où elles sont
données : les démontrer demanderait une théorie de la mesure. Ce qui se démontre, ce sont
les relations entre elles, et leur comportement par agrandissement.
Énoncés et démonstrations en français : voir GrandeursEtMesures.tex.

\subsection{Définitions}

\noindent
Toutes les définitions du fichier sont posées ici, avant les démonstrations.

\begin{definition}
Le \emph{périmètre du cercle} de rayon $r$ vaut $\mathcal{P}(r) = 2\pi r$.
\end{definition}
\source{https://github.com/Commutator-IO/learn-lean/blob/main/courses/01-college/04-grandeurs-et-mesures/GrandeursEtMesures.lean\#L20}{GrandeursEtMesures.lean\#L20}

\begin{definition}
L'\emph{aire du disque} de rayon $r$ vaut $\mathcal{A}(r) = \pi r^2$.
\end{definition}
\source{https://github.com/Commutator-IO/learn-lean/blob/main/courses/01-college/04-grandeurs-et-mesures/GrandeursEtMesures.lean\#L23}{GrandeursEtMesures.lean\#L23}

\begin{definition}
L'\emph{aire d'un rectangle} de côtés $L$ et $\ell$ vaut $L\ell$.
\end{definition}
\source{https://github.com/Commutator-IO/learn-lean/blob/main/courses/01-college/04-grandeurs-et-mesures/GrandeursEtMesures.lean\#L26}{GrandeursEtMesures.lean\#L26}

\begin{definition}
Le \emph{périmètre d'un rectangle} de côtés $L$ et $\ell$ vaut $2(L + \ell)$.
\end{definition}
\source{https://github.com/Commutator-IO/learn-lean/blob/main/courses/01-college/04-grandeurs-et-mesures/GrandeursEtMesures.lean\#L29}{GrandeursEtMesures.lean\#L29}

\begin{definition}
L'\emph{aire d'un triangle} de base $b$ et de hauteur $h$ vaut $\dfrac{bh}{2}$.
\end{definition}
\source{https://github.com/Commutator-IO/learn-lean/blob/main/courses/01-college/04-grandeurs-et-mesures/GrandeursEtMesures.lean\#L32}{GrandeursEtMesures.lean\#L32}

\begin{definition}
L'\emph{aire d'un parallélogramme} de base $b$ et de hauteur $h$ vaut $bh$.
\end{definition}
\source{https://github.com/Commutator-IO/learn-lean/blob/main/courses/01-college/04-grandeurs-et-mesures/GrandeursEtMesures.lean\#L35}{GrandeursEtMesures.lean\#L35}

\begin{definition}
L'\emph{aire d'un trapèze} de bases $B$ et $b$ et de hauteur $h$ vaut
$\dfrac{(B + b)h}{2}$.
\end{definition}
\source{https://github.com/Commutator-IO/learn-lean/blob/main/courses/01-college/04-grandeurs-et-mesures/GrandeursEtMesures.lean\#L38}{GrandeursEtMesures.lean\#L38}

\begin{definition}
Le \emph{volume d'un pavé droit} de dimensions $L$, $\ell$, $H$ vaut $L\ell H$.
\end{definition}
\source{https://github.com/Commutator-IO/learn-lean/blob/main/courses/01-college/04-grandeurs-et-mesures/GrandeursEtMesures.lean\#L41}{GrandeursEtMesures.lean\#L41}

\begin{definition}
Le \emph{volume d'un prisme} ou d'un cylindre est le produit de l'aire de sa base par
sa hauteur.
\end{definition}
\source{https://github.com/Commutator-IO/learn-lean/blob/main/courses/01-college/04-grandeurs-et-mesures/GrandeursEtMesures.lean\#L44}{GrandeursEtMesures.lean\#L44}

\begin{definition}
Le \emph{volume d'une pyramide} ou d'un cône est le tiers de celui du prisme de mêmes
base et hauteur.
\end{definition}
\source{https://github.com/Commutator-IO/learn-lean/blob/main/courses/01-college/04-grandeurs-et-mesures/GrandeursEtMesures.lean\#L47}{GrandeursEtMesures.lean\#L47}

\begin{definition}
Le \emph{volume de la boule} de rayon $r$ vaut $\dfrac{4}{3}\pi r^3$.
\end{definition}
\source{https://github.com/Commutator-IO/learn-lean/blob/main/courses/01-college/04-grandeurs-et-mesures/GrandeursEtMesures.lean\#L50}{GrandeursEtMesures.lean\#L50}

\subsection{Périmètre du cercle et aire du disque}

\begin{theoreme}
Le périmètre du cercle est proportionnel au rayon, de coefficient $2\pi$ :
$\mathcal{P}(r) = 2\pi r$.
\end{theoreme}

\begin{proof}
C'est la définition même du périmètre, lue comme une proportionnalité : le rapport
$\mathcal{P}(r)/r$ vaut $2\pi$ quel que soit $r$ non nul.
\end{proof}
\source{https://github.com/Commutator-IO/learn-lean/blob/main/courses/01-college/04-grandeurs-et-mesures/GrandeursEtMesures.lean\#L55}{GrandeursEtMesures.lean\#L55}

\begin{theoreme}
L'aire du disque n'est pas proportionnelle au rayon mais à son carré : doubler le
rayon quadruple l'aire,
\[ \mathcal{A}(2r) = 4\,\mathcal{A}(r) . \]
\end{theoreme}

\begin{proof}
On remplace $r$ par $2r$ dans la formule :
\[ \mathcal{A}(2r) = \pi (2r)^2 = 4\pi r^2 = 4\,\mathcal{A}(r) . \]
\end{proof}
\source{https://github.com/Commutator-IO/learn-lean/blob/main/courses/01-college/04-grandeurs-et-mesures/GrandeursEtMesures.lean\#L59}{GrandeursEtMesures.lean\#L59}

\subsection{Aires du rectangle, du triangle, du parallélogramme et du trapèze}

\begin{theoreme}
Le triangle a pour aire la moitié de celle du parallélogramme de mêmes base et
hauteur : $\dfrac{bh}{2} = \dfrac{1}{2}(bh)$.
\end{theoreme}

\begin{proof}
Les deux formules coïncident après division par deux. Géométriquement, deux copies du
triangle assemblées forment le parallélogramme.
\end{proof}
\source{https://github.com/Commutator-IO/learn-lean/blob/main/courses/01-college/04-grandeurs-et-mesures/GrandeursEtMesures.lean\#L67}{GrandeursEtMesures.lean\#L67}

\begin{theoreme}
Un parallélogramme et un rectangle de mêmes base et hauteur ont la même aire, $bh$.
\end{theoreme}

\begin{proof}
Les deux formules sont la même. C'est le découpage-recollement classique : on détache le
triangle qui dépasse d'un côté pour le rapporter de l'autre.
\end{proof}
\source{https://github.com/Commutator-IO/learn-lean/blob/main/courses/01-college/04-grandeurs-et-mesures/GrandeursEtMesures.lean\#L72}{GrandeursEtMesures.lean\#L72}

\begin{theoreme}
Un trapèze dont les deux bases sont égales est un parallélogramme :
$\dfrac{(b + b)h}{2} = bh$.
\end{theoreme}

\begin{proof}
On calcule :
\[ \frac{(b+b)h}{2} = \frac{2bh}{2} = bh . \]
\end{proof}
\source{https://github.com/Commutator-IO/learn-lean/blob/main/courses/01-college/04-grandeurs-et-mesures/GrandeursEtMesures.lean\#L77}{GrandeursEtMesures.lean\#L77}

\begin{theoreme}
Le trapèze a l'aire du rectangle construit sur la moyenne de ses bases :
\[ \frac{(B + b)h}{2} = \frac{B + b}{2} \times h . \]
\end{theoreme}

\begin{proof}
Les deux écritures ne diffèrent que par le placement de la division par deux.
\end{proof}
\source{https://github.com/Commutator-IO/learn-lean/blob/main/courses/01-college/04-grandeurs-et-mesures/GrandeursEtMesures.lean\#L82}{GrandeursEtMesures.lean\#L82}

\subsection{Deux figures de même aire peuvent avoir des périmètres différents}

\begin{theoreme}
Deux rectangles de même aire peuvent avoir des périmètres différents : $1 \times 12$ et
$3 \times 4$ ont pour aire $12$, et pour périmètres $26$ et $14$.
\end{theoreme}

\begin{proof}
On calcule les quatre nombres : $1 \times 12 = 3 \times 4 = 12$, tandis que
$2(1 + 12) = 26$ et $2(3 + 4) = 14$. L'aire ne détermine donc pas le périmètre.
\end{proof}
\source{https://github.com/Commutator-IO/learn-lean/blob/main/courses/01-college/04-grandeurs-et-mesures/GrandeursEtMesures.lean\#L90}{GrandeursEtMesures.lean\#L90}

\begin{theoreme}
Réciproquement, deux rectangles de même périmètre peuvent avoir des aires différentes :
$1 \times 5$ et $3 \times 3$ ont pour périmètre $12$, et pour aires $5$ et $9$.
\end{theoreme}

\begin{proof}
Mêmes calculs : $2(1+5) = 2(3+3) = 12$, alors que $1 \times 5 = 5$ et $3 \times 3 = 9$.
\end{proof}
\source{https://github.com/Commutator-IO/learn-lean/blob/main/courses/01-college/04-grandeurs-et-mesures/GrandeursEtMesures.lean\#L97}{GrandeursEtMesures.lean\#L97}

\subsection{Volumes}

\begin{theoreme}
La pyramide a pour volume le tiers de celui du prisme de mêmes base et hauteur.
\end{theoreme}

\begin{proof}
Les deux formules ne diffèrent que par le facteur $\tfrac13$, qui est celui de la
définition. Le collège l'admet ; le démontrer demande un calcul intégral.
\end{proof}
\source{https://github.com/Commutator-IO/learn-lean/blob/main/courses/01-college/04-grandeurs-et-mesures/GrandeursEtMesures.lean\#L105}{GrandeursEtMesures.lean\#L105}

\begin{theoreme}
Le pavé droit est le prisme de base rectangulaire : $L\ell H$ est l'aire de base
$L\ell$ multipliée par la hauteur $H$.
\end{theoreme}

\begin{proof}
Les deux expressions sont identiques après regroupement des facteurs.
\end{proof}
\source{https://github.com/Commutator-IO/learn-lean/blob/main/courses/01-college/04-grandeurs-et-mesures/GrandeursEtMesures.lean\#L110}{GrandeursEtMesures.lean\#L110}

\subsection{Agrandissement-réduction de rapport \texttt{k} : longueurs \texttt{\(\times\) k}, aires \texttt{\(\times\) k²}, volumes \texttt{\(\times\) k³}}

\begin{theoreme}
Un agrandissement de rapport $k$ multiplie les longueurs par $k$ :
\[ \mathcal{P}(kL, k\ell) = k\,\mathcal{P}(L, \ell) . \]
\end{theoreme}

\begin{proof}
On factorise :
\[ 2(kL + k\ell) = k \times 2(L + \ell) . \]
\end{proof}
\source{https://github.com/Commutator-IO/learn-lean/blob/main/courses/01-college/04-grandeurs-et-mesures/GrandeursEtMesures.lean\#L116}{GrandeursEtMesures.lean\#L116}

\begin{theoreme}
Il multiplie les aires par $k^2$ :
\[ \mathcal{A}(kL, k\ell) = k^2\,\mathcal{A}(L, \ell) . \]
\end{theoreme}

\begin{proof}
Chacune des deux dimensions est multipliée par $k$ :
\[ (kL)(k\ell) = k^2 L\ell . \]
\end{proof}
\source{https://github.com/Commutator-IO/learn-lean/blob/main/courses/01-college/04-grandeurs-et-mesures/GrandeursEtMesures.lean\#L122}{GrandeursEtMesures.lean\#L122}

\begin{theoreme}
Même chose pour le disque, où le rapport porte sur le rayon :
$\mathcal{A}(kr) = k^2\,\mathcal{A}(r)$.
\end{theoreme}

\begin{proof}
$\pi (kr)^2 = k^2 \pi r^2$.
\end{proof}
\source{https://github.com/Commutator-IO/learn-lean/blob/main/courses/01-college/04-grandeurs-et-mesures/GrandeursEtMesures.lean\#L128}{GrandeursEtMesures.lean\#L128}

\begin{theoreme}
Il multiplie les volumes par $k^3$ :
\[ \mathcal{V}(kL, k\ell, kH) = k^3\,\mathcal{V}(L, \ell, H) . \]
\end{theoreme}

\begin{proof}
Les trois dimensions sont multipliées par $k$ :
\[ (kL)(k\ell)(kH) = k^3 L\ell H . \]
C'est la raison pour laquelle un modèle réduit au dixième pèse mille fois moins.
\end{proof}
\source{https://github.com/Commutator-IO/learn-lean/blob/main/courses/01-college/04-grandeurs-et-mesures/GrandeursEtMesures.lean\#L133}{GrandeursEtMesures.lean\#L133}

\begin{theoreme}
Même chose pour la boule : $\mathcal{V}(kr) = k^3\,\mathcal{V}(r)$.
\end{theoreme}

\begin{proof}
$\tfrac43 \pi (kr)^3 = k^3 \times \tfrac43 \pi r^3$.
\end{proof}
\source{https://github.com/Commutator-IO/learn-lean/blob/main/courses/01-college/04-grandeurs-et-mesures/GrandeursEtMesures.lean\#L139}{GrandeursEtMesures.lean\#L139}

\vfill
\noindent\rule{0.35\linewidth}{0.4pt}\par
\noindent{\footnotesize Leonardo de Moura et Sebastian Ullrich,
\emph{The Lean~4 Theorem Prover and Programming Language},
\emph{Automated Deduction — CADE~28}, Springer, 2021, p.~625--635,
\href{https://doi.org/10.1007/978-3-030-79876-5_37}{doi:10.1007/978-3-030-79876-5\_37}.\par}

\end{document}
